% from ega4-1-fr.tex
\oldpage[IV-1]{224}

quelconque) et par image directe (par un morphisme de présentation finie), et
d'avoir des propriétés topologiques maniables dans les préschémas localement
noethériens. Les mêmes techniques permettent même souvent de se ramener à des
anneaux noethériens plus particuliers, par exemple les
$\mathbf Z$-\emph{algèbres de type fini}, et c'est là que les propriétés des
anneaux «~excellents~» (étudiées au §~7) interviennent de façon décisive.
Indépendamment d'ailleurs de la question de l'élimination des hypothèses
noethériennes, les techniques des §§~8 à 10, de nature très élémentaire, sont
d'usage constant dans presque toutes les applications.

\setcounter{section}{0}

\section{Conditions de finitude relatives. Ensembles constructibles dans les préschémas}
\label{section:IV.1-fr}

Nous allons, dans ce paragraphe, reprendre, en les complétant, l'exposé des
«~conditions de finitude~» pour un morphisme de préschémas $f:X\to Y$, donné
dans (I, 6.3 et 6.6). Il y a essentiellement deux notions de «~finitude~» de
nature \emph{globale} sur $X$, celle de morphisme \emph{quasi-compact}
(définie dans (I, 6.6.1)) et celle de morphisme \emph{quasi-séparé}~; il y a
d'autre part deux notions de «~finitude~» de nature \emph{locale} sur $X$,
celle de morphisme \emph{localement de type fini} (définie dans (I, 6.6.2))
et celle de morphisme \emph{localement de présentation finie}. En combinant
ces notions locales aux notions globales précédentes, on obtient les notions
de morphisme \emph{de type fini} (définie dans (I, 6.3.1)) et de morphisme
\emph{de présentation finie}. Pour la commodité du lecteur, nous donnerons de
nouveau, dans ce paragraphe, les principales propriétés énoncées dans (I, 6.3
et 6.6), en renvoyant bien entendu à ces numéros du chapitre I pour leurs
démonstrations.

Aux n\textsuperscript{os}~(1.8) et (1.9), nous complétons, dans le cadre des
préschémas, et en faisant usage des notions de finitude précédentes, les
résultats sur les ensembles constructibles donnés dans
(0\textsubscript{III}, §~9).

\subsection{Morphismes quasi-compacts}
\label{subsection:IV.1.1-fr}

\begin{definition}[1.1.1]
\label{IV.1.1.1-fr}
On dit qu'un morphisme de préschémas $f:X\to Y$ est \emph{quasi-compact} si
l'application continue $f$ de l'espace topologique $X$ dans l'espace
topologique $Y$ est quasi-compacte (0\textsubscript{III}, 9.1.1), autrement
dit si l'image réciproque $f^{-1}(U)$ de tout ouvert quasi-compact $U$ de $Y$
est quasi-compacte (cf. (I, 6.6.1)).
\end{definition}

Si $\mathfrak B$ est une base de la topologie de $Y$ formée d'ouverts
affines, pour que $f$ soit quasi-compact, il faut et il suffit que pour tout
$V\in\mathfrak B$, $f^{-1}(V)$ soit \emph{réunion finie d'ouverts affines}.
Par exemple, si $Y$ est affine et $X$ quasi-compact, \emph{tout} morphisme
$f:X\to Y$ est quasi-compact (I, 6.6.1).

Si $f:X\to Y$ est un morphisme quasi-compact, il est clair que pour tout
ouvert $V$ de $Y$, la restriction de $f$ à $f^{-1}(V)$ est un morphisme
quasi-compact $f^{-1}(V)\to V$. Inversement, si $(U_\alpha)$ est un
recouvrement ouvert de $Y$ et $f:X\to Y$ un morphisme tel que les restrictions
$f^{-1}(U_\alpha)\to U_\alpha$ soient quasi-compactes, alors $f$ est
quasi-compact. Par

